\subsection{Support Vector Machines (SVM)}
\label{subsubsec:eda_svm}

Las Máquinas de Vectores de Soporte (\textit{Support Vector Machines}, SVM) son un conjunto de algoritmos de aprendizaje supervisado desarrollados por Vladimir Vapnik y su equipo en AT\&T Bell Laboratories durante la década de 1990 ~\parencite{1995svm}. Originalmente diseñadas para clasificación binaria, las SVM se han convertido en uno de los métodos más robustos y teóricamente fundamentados del aprendizaje automático.

\subsubsection*{Fundamento y funcionamiento}

El objetivo fundamental de una SVM es encontrar un hiperplano que separe las clases (en nuestro caso dos, voces reales y voces generadas por IA) con el máximo margen posible. La clave de reside en que no cualquier línea divisoria es válida: el algoritmo busca aquella que maximiza la distancia entre la línea y los puntos más cercanos de cada clase. Estos puntos críticos, que ``sostienen'' el margen, reciben el nombre de \textbf{vectores de soporte} (\textit{support vectors}), de donde el algoritmo toma su nombre.

El proceso de entrenamiento de una SVM puede entenderse en los siguientes pasos:

\begin{enumerate}
    \item \textbf{Identificación de vectores de soporte}: El algoritmo examina todos los puntos de entrenamiento (instancias) e identifica aquellas que son más difíciles de clasificar, es decir, las que están más cerca de la frontera entre clases.
    
    \item \textbf{Maximización del margen}: A partir de estos vectores de soporte, la SVM calcula el hiperplano que maximiza la distancia a ambos lados. Este margen máximo proporciona la mejor capacidad de generalización: cuanto más ancho sea el margen, menor será la probabilidad de error al clasificar nuevas instancias.
    
    \item \textbf{Clasificación de nuevas muestras}: Una vez entrenado el modelo, para clasificar una nueva instancia, simplemente se calcula a qué lado del hiperplano se sitúa.
\end{enumerate}

\subsubsection*{El truco del kernel: Separando lo inseparable}

Uno de los mayores desafíos en la clasificación es que las instancias rara vez son linealmente separables en el espacio original de características. Es decir, no existe una línea recta (o hiperplano) que pueda separarlas perfectamente. Para abordar esta limitación, las SVM introducen el concepto de \textit{kernel} o núcleo.

El \textit{truco del kernel} funciona de la siguiente manera conceptual:

\begin{itemize}
    \item Los datos originales (no separables linealmente) se transforman implícitamente a un espacio de mayor dimensionalidad mediante una función matemática.
    \item En este nuevo espacio, es posible que los datos sí sean linealmente separables.
    \item La ventaja del kernel es que realiza esta transformación sin necesidad de calcular explícitamente las coordenadas en el espacio de alta dimensión, lo que sería computacionalmente costoso.
\end{itemize}

Los kernels más comúnmente utilizados son:

\begin{itemize}
    \item \textbf{Kernel lineal}: Asume que los datos son aproximadamente separables por un hiperplano. Útil cuando el número de características es muy grande.
    
    \item \textbf{Kernel polinomial}: Crea fronteras de decisión curvas mediante combinaciones polinómicas de las características originales.
    
    \item \textbf{Kernel RBF (\textit{Radial Basis Function})}: Es el más utilizado en muchos problemas. Proyecta los datos a un espacio de dimensionalidad infinita, permitiendo fronteras de decisión altamente complejas y no lineales. Se basa en la similitud entre muestras medida por distancia euclidiana:

    \[
    K(\mathbf{x}_i, \mathbf{x}_j) = \exp(-\gamma \|\mathbf{x}_i - \mathbf{x}_j\|^2)
    \]

    donde \(\gamma\) controla la influencia de cada muestra: un valor pequeño considera vecindarios amplios (frontera más suave), mientras que un valor grande se enfoca en vecindarios reducidos (frontera más ajustada a los datos).
\end{itemize}

\subsubsection*{Hiperparámetros}

Veamos cuáles son los hiperparámetros más influyentes en los SVM:

\begin{itemize}
    \item \textbf{C}: Controla el equilibrio entre tener un margen amplio y clasificar correctamente los puntos de entrenamiento.
    
    \item \textbf{kernel}: Es la función que transforma los datos a un espacio de mayor dimensión donde puedan ser separables linealmente.
    
    \item \textbf{gamma \(\gamma\)}: Específico de kernels no lineales (RBF, sigmoide o polinomial). Define cuánta influencia tiene un solo punto de entrenamiento.
    
    \item \textbf{probability}: Este parámetro $booleano$ determina si el modelo SVM debe calibrarse para proporcionar estimaciones de probabilidad junto con las predicciones de clase.
\end{itemize}

\subsubsection*{Ventajas específicas para detección de deepfakes}

Repasemos algunas de las ventajas que presentan los SVM para la detección de $deepfakes$:

\begin{itemize}
    \item \textbf{Fundamento teórico sólido}: A diferencia de otros métodos heurísticos, SVM está basada en la teoría de optimización convexa, lo que garantiza que la solución encontrada es óptima (máximo margen global).
    
    \item \textbf{Eficaz en espacios de alta dimensión}: Las características de audio pueden generar vectores de gran dimensionalidad (decenas o cientos de coeficientes). SVM maneja bien esta situación sin colapsar por la maldición de la dimensionalidad.
    
    \item \textbf{Robustez ante overfitting}: La maximización del margen proporciona una regularización natural que mejora la generalización, crucial cuando se trabaja con datasets limitados de voces $deepfake$.
    
    \item \textbf{Versatilidad mediante kernels}: La capacidad de elegir diferentes kernels permite adaptar el modelo a la estructura no lineal de los datos de audio sin necesidad de transformaciones manuales.
\end{itemize}